%%%%%%%%%%%%%%%%%%%%%%%%%%%%%%%%%%%%%%%%%%%%%%%%%%%%%%%%%%
%%
%% MARK: Title
%%
%%%%%%%%%%%%%%%%%%%%%%%%%%%%%%%%%%%%%%%%%%%%%%%%%%%%%%%%%%

\chapter*{\S B6. Basic $1$-Forms on Principal Fiber Bundles}
\setcounter{footnote}{0}

%%%%%%%%%%%%%%%%%%%%%%%%%%%%%%%%%%%%%%%%%%%%%%%%%%%%%%%%%%
%%
%% MARK: Abstract
%%
%%%%%%%%%%%%%%%%%%%%%%%%%%%%%%%%%%%%%%%%%%%%%%%%%%%%%%%%%%

\vspace{-1ex}\begin{abstract} % <- PIZ adjust
  We give a criterion for basic $1$-forms on general principal bundles.
\end{abstract}

%%%%%%%%%%%%%%%%%%%%%%%%%%%%%%%%%%%%%%%%%%%%%%%%%%%%%%%%%%
%%
%% MARK: Content
%%
%%%%%%%%%%%%%%%%%%%%%%%%%%%%%%%%%%%%%%%%%%%%%%%%%%%%%%%%%%

\vspace{-1ex}This is a specialization of the general criterion of paragraph 6.38 in the textbook, % <- PIZ adjust
for basic $1$-forms on principal fiber bundles \cite[\textsection\,8.11]{PIZ13}.

\begin{Article}{Exercise.}{}
  Let $\pi : X \to B$ be a principal fibration with structure group $G$.
  Let $\alpha$ be a $1$-form on $X$.
  Prove that $\alpha$ is {\em basic},
  that is, is the pullback $\pi^*(\beta)$ of a $1$-form
  on $B$,
  if and only if:
  \begin{enumerate}
    \item $\alpha$  vanishes on ``vertical plots'', that is, the preimages
    by $\pi$ of the points of $B$.
    In short, $\alpha \restriction \pi^{-1}(b) = 0$,
    for all $b \in B$.
    \item $\alpha$  is invariant by the structure group,
    that is,
    $g_X^*(\alpha) = \alpha$,
    for all $g \in G$,
    where $g_X$ denotes the action of $g$ on $X$.
  \end{enumerate}
  Give a simple example showing that this criterion is then no longer valid even for a
  $2$-form.
\end{Article}

\begin{proof}
  Assume that $\alpha = \pi^*(\beta)$.
  Let $P : U \to X$ be a vertical plot,
  then $\alpha(P) = [\pi^*(\beta)](P) = \beta(\pi \circ P)$,
  but $\pi \circ P = \Const$,
  thus $\beta(\pi \circ P) = 0$,
  that is,
  $\alpha(P) = 0$,
  see \cite[Exercise 96]{PIZ13}.
  Now,
  let $g \in G$,
  then $g_X^*(\alpha) = g_X^*(\pi^*(\beta)) = (\pi \circ g_X)^*(\beta)$,
  but $\pi \circ g_X = \pi$,
  thus $g_X^*(\alpha) = \alpha$.
  Therefore,
  the two conditions above are satisfied.
  
  Conversely,
  assume that the two conditions above are satisfied,
  and let $P : U \to X$ and $P' : U \to X$ be two plots such that $\pi \circ P = \pi \circ P'$.
  Since $\pi$ is a principal diffeological fibration,
  there exists, around every point of $U$,
  a subdomain $V \subset U$ and a plot $r \mapsto \gamma(r)$ in $G$,
  defined on $V$,
  such that
  $P'(r) = \gamma(r)_X(P(r))$,
  for all $r \in V$ \cite[8.13, Note 1]{PIZ13}.
  Hence,
  \vspace{-.33\baselineskip} % <- PIZ adjust
  $$%
    \alpha(P') = \alpha[r \mapsto \gamma(r)_X(P(r))].
    $$%
  Now we need the following formula,
  which is a generalization of
  \cite[8.37 $(\clubsuit)$]{PIZ13}.
  Let $\Rm(x) : G \to X$ be
  the orbit map $\Rm(x)(g) = g_X(x)$, let $\gamma$ be a plot in $G$
  defined on some domain $U$ and let $P$ be a plot
  in $X$ defined on some domain $V$,
  then
  \begin{multline*}
    \alpha[(r,s) \mapsto \gamma(r)_X(P(s))]_{r \choose s}\Vector{\delta r \\ \delta s} = \alpha[(r,s) \mapsto \gamma(r)_X(P(s))]_{r \choose s}\Vector{\delta r \\ 0} \\
    + \alpha[(r,s) \mapsto \gamma(r)_X(P(s))]_{r \choose s}\Vector{0 \\ \delta s}  = \alpha[r \mapsto \gamma(r)_X(P(s))]_r (\delta r) \\
    + \alpha[s \mapsto \gamma(r)_X(P(s))]_s(\delta s) =  [\Rm(P(s))^*(\alpha)](\gamma)_r(\delta r) +  [\gamma(r)_X^*(\alpha)](P)_s(\delta s).
  \end{multline*}
  Now,
  since $\Rm(P(s))$ is the orbit map of the point $x = P(s)$,
  $$
    [\Rm(P(s))^*(\alpha)](\gamma) = [\Rm(P(s))^*(\alpha \restriction \pi^{-1}(b))](\gamma),
    $$
  with $b = \pi(x)$. But, according to the hypothesis, $\alpha \restriction \pi^{-1}(b) = 0$, thus
  $$
    [\Rm(P(s))^*(\alpha)](\gamma)_r(\delta r) = 0.
    $$
  Next,
  since,
  by assumption,
  $\alpha$ is invariant by the action of $G$,
  $[\gamma(r)_X^*(\alpha)](P) = \alpha(P)$.
  Therefore,
  $$
    \alpha[(r,s) \mapsto \gamma(r)_X(P(s))]_{r \choose s}\Vector{\delta r \\ \delta s}
    = \alpha(P)_r(\delta r).
    $$
  This implies in particular $\alpha(P') = \alpha[r \mapsto \gamma(r)_X(P(r))] =
  \alpha(P)$.
  Hence,
  according to the criterion \cite[6.38]{PIZ13},
  there exists a $1$-form
  $\beta$ on $B$ such that $\alpha = \pi^*(\beta)$.
  
  For the second question, consider the $2$-form $\omega = dx \wedge dy$ on the plane $\RR^2$.
  We can consider the first projection $\Proj_1 : (x,y) \mapsto x$.
  It is a
  principal fibration with group $(\RR,+)$ for the action
  $t_{\RR^2}(x,y) = (x, y+t)$.
   The $2$-form $\omega$ is,
  in particular,
  invariant by this action.
  Moreover,
  the restriction of $\omega$ on a fiber of $\Proj_1$ is a linear $2$-form
  on $\RR$,
  thus vanishes.
  This example satisfies the two conditions above,
  but since $\omega \neq 0$, $\omega$ is not the pullback by $\Proj_1$ of a $2$-form on $\RR$
  (which is necessarily $0$).
\end{proof}
